\section{Faithful cyclotomic invariants and the actual etale functor}
\label{sec:faithful-etale-v057}

\status{This chapter proves the original faithful-invariants statement,
constructs a fixed basis on the prescribed integral Witt tensor lattice,
constructs the Frobenius change of basis into the original coefficient
ring, and returns the full original etale functor, including its maps on
all lattice morphisms. The qualitative perfectoid finite-extension
theorem is used as a precisely stated classical foundation for the first
step. All later averaging and conjugacy arguments are written out.
The original good layer and period records remain supplied data.
Herr cohomology, the quantitative comparison and the common global
denominator family remain separate.}

The problem is to make the action of the cyclotomic kernel disappear
from a basis without changing either the lattice or its coefficient
embedding. There are two steps. A finite averaging argument followed by
an integral limit produces a fixed basis. Its Frobenius matrix initially
has entries in a larger fixed ring. A second, nonlinear change of basis
puts those entries in the specified smaller ring. The previously proved
IFD25 concerns descent of a matrix already satisfying an equation; it
does not by itself construct this second change of basis. Below we give
that construction before using IFD25 to descend morphisms.

Fix a prime $p$; the argument includes $p=2$. Let
$\overline{\mathbb Q}_p\subset C$ be the original algebraic closure
and its completion. Put
\[
 S=\mathcal O_C^\flat=\varprojlim_{x\mapsto x^p}\mathcal O_C/p,
 \qquad F=\operatorname{Frac}(S),\qquad W_0=W(F).
                                                        \tag{FET.1}
\]
For a finite coefficient extension $LL/\mathbb Q_p$ put
$O=\mathcal O_{LL}$ and $\mathcal W=O\otimes_{\mathbb Z_p}W_0$.
For a continuous $O$-lattice $T$ with $G_K$-action put
$M_T=\mathcal W\otimes_O T$ and $H=G_K\cap\ker\chi$.
The notation $T$ refers to the integral lattice itself. Its scalar
extension to $LL$ is not substituted for it.

The original records $P,AI$ prescribe the actions on $C,S,W_0$ and
$\mathcal W$. Their identity with the actions just described is part of
the proof, rather than a change in the input data. Later, on a good
layer, $E=R.\mathrm{EE}$ and $j=I.\mathrm{ioE}:E\to\mathcal W$
will denote the original coefficient ring and map.

\subsection{Why the tilt separates the original finite extensions}

% BEGIN FET1
\paragraph{FET1: the prescribed actions are the canonical actions.}
\label{claim:FET1}
For every original $P,AI$, the action $AI.\mathrm{actC}$ is the
unique isometric extension of the given action on
$\overline{\mathbb Q}_p$. The action $AI.\mathrm{actT}$ is its
coordinatewise action on the inverse limit in (FET.1). Moreover,
$\ker(P.\chi_Q)$ fixes every $p$-power root of unity, and is exactly
the usual cyclotomic kernel.

\begin{proof}
Two isometries of $C$ which agree on the dense subfield
$\overline{\mathbb Q}_p$ agree on a convergent sequence approaching
each point of $C$. The fields \texttt{actC\_ext} and
\texttt{actC\_norm} therefore determine the first action uniquely.
For $x\in S$, every coordinate of $AI.\mathrm{actT}(g)x$ is
specified by \texttt{actT\_coeff}: lift the corresponding residue
coordinate of $x$ to $\mathcal O_C$, apply the already determined
isometry, and reduce modulo $p$. Equality of all coordinates proves
the second assertion.

For a primitive $p^m$-th root $\zeta$, the field
\texttt{chiQ\_cyc} gives an integer $a$ whose distance from
$P.\chi_Q(g)$ is at most $p^{-m}$ and for which $g\zeta=\zeta^a$.
The ordinary cyclotomic exponent has the same residue modulo $p^m$.
For all $m$ these residues determine one element of $\mathbb Z_p$,
so it equals $P.\chi_Q(g)$ in $\mathbb Q_p$.
This also follows from the original integral-character result WPR2.
In particular a value $1$ fixes every root, and the converse holds by
uniqueness of the residues. No alternate action or character has been
introduced.
\end{proof}
% END FET1

% BEGIN FET2
\paragraph{FET2: the completed cyclotomic extension of a local field.}
\label{claim:FET2}
For a finite extension $L/\mathbb Q_p$ in the original algebraic
closure, let
\[
 L_\infty=L(\mu_{p^\infty}),\qquad
 E_L=\overline{L_\infty}^{\,C}.
\]
Then $E_L$ is a perfectoid field. In particular its ring of integers
is complete for the $p$-adic topology and Frobenius on
$\mathcal O_{E_L}/p$ is onto.

\begin{proof}
Write $E_0=\widehat{\mathbb Q_p(\mu_{p^\infty})}$.
First one can see its perfectoid property directly. The integral
ring at level $m$ is $\mathbb Z_p[\zeta_{p^m}]$; the usual Eisenstein
polynomial for $\zeta_{p^m}-1$ proves this assertion, including the
initial trivial levels when $p=2$. Modulo $p$, a polynomial in
$\zeta_{p^m}$ has a $p$-th root at level $m+1$: reduce its
coefficients to $\mathbb F_p$ and replace $\zeta_{p^m}$ by
$\zeta_{p^{m+1}}$. Every integral element of the completion is
congruent modulo $p$ to an integral element at a finite level, by
density. Thus Frobenius on $\mathcal O_{E_0}/p$ is onto.
The valuations of $\zeta_{p^m}-1$, namely
$1/(p^{m-1}(p-1))$, show that the value group is nondiscrete.
Together with completeness these are the perfectoid-field conditions.

We now use the classical finite-extension theorem for perfectoid
fields: every finite extension of a perfectoid field, with its natural
complete valuation topology, is perfectoid. This is Scholze,
\emph{Perfectoid spaces}, Theorem~3.7(i); this named foundation is
the only deep input in the present separation argument.
The compositum $LE_0$ inside $C$ has finite degree over $E_0$ and
is complete, hence closed in $C$. It contains $L_\infty$.
Conversely $E_L$ contains both $L$ and $E_0$, hence $LE_0$.
Thus $E_L=LE_0$, so the theorem applies to the actual embedded field.
Equivalently, Bhatt--Morrow--Scholze, \emph{Integral $p$-adic Hodge
theory}, Example~3.6, states this integral perfectoid example directly.
There is no Galois or unramified assumption on $L$.
\end{proof}
% END FET2

% BEGIN FET3
\paragraph{FET3: recovering integral elements from their tilt coordinates.}
\label{claim:FET3}
Let $E\subset C$ be a closed perfectoid subfield containing
$\mathbb Q_p$. Its integral tilt
$S_E=\varprojlim_{x\mapsto x^p}\mathcal O_E/p$ embeds in $S$.
For $x=(\overline x_m)\in S_E$ and arbitrary lifts
$x_m\in\mathcal O_E$, the limit
\[
                     x^\sharp=\lim_m x_m^{p^m}             \tag{FET.2}
\]
exists, is independent of the lifts, and reduces to
$\overline x_0$ modulo $p$. Every $a\in\mathcal O_E$ admits an
expansion
\[
                      a=\sum_{n\geq0}p^n x_n^\sharp,
                      \qquad x_n\in S_E.                 \tag{FET.3}
\]
If an isometric automorphism of $C$ fixes every image of $S_E$ in
$S$, it fixes $E$ pointwise.

\begin{proof}
The injection $\mathcal O_E/p\hookrightarrow\mathcal O_C/p$
follows from $p\mathcal O_C\cap\mathcal O_E=p\mathcal O_E$:
divide an element of the intersection by $p$ in $E$ and use its
valuation. Passing to inverse limits is therefore injective.
If $u-v\in p^a\mathcal O_E$, with $a\geq1$, then the binomial
formula gives $u^p-v^p\in p^{a+1}\mathcal O_E$. The terms with
$1\leq i<p$ have a factor $p$, while $pa\geq a+1$ handles the
last term, also for $p=2$. Consequently compatible residue lifts
satisfy
\[
 x_{m+1}^{p^{m+1}}-x_m^{p^m}\in p^{m+1}\mathcal O_E.
\]
Changing a lift has the same estimate. Completeness proves (FET.2)
and its stated residue property.

Surjectivity of Frobenius permits successive choices of roots of
any residue class, so the zeroth projection
$S_E\to\mathcal O_E/p$ is onto. Given $a_0=a$, choose $x_0$
whose sharp has the same residue, and set
$a_1=(a_0-x_0^\sharp)/p\in\mathcal O_E$. Repeat.
The error after $n$ terms is $p^n a_n$, proving (FET.3).
Formula (FET.2) is equivariant for isometric automorphisms, even
when the automorphism is not assumed to preserve $E$ in advance.
Fixing all the $x_n$ therefore fixes every sum in (FET.3).
Finally every element of $E$ becomes integral after multiplication
by a sufficiently large power of $p$. This proves the last assertion.
The sharp map has not been treated as additive.
\end{proof}
% END FET3

% BEGIN FET4
\paragraph{FET4: full original faithful cyclotomic invariants.}
\label{claim:FET4}
For every finite $L/\mathbb Q_p$ and every
$g\in G_{\mathbb Q_p}$, if $AI.\mathrm{actT}(g)$ fixes every
element of $S$ fixed by $G_L\cap\ker\chi$, then $g\in G_L$.
In particular the full original predicate \texttt{FaithH AI} holds.

\begin{proof}
Every $h\in G_L\cap\ker\chi$ fixes $L_\infty$ pointwise.
Its isometric extension fixes the closure $E_L$ by FET1, hence
fixes every coordinate of $S_{E_L}\subset S$.
The assumption on $g$ therefore fixes all of $S_{E_L}$.
FET2 and FET3 imply that $g$ fixes $E_L$, and hence $L$.
This is precisely membership in $G_L$.
The original statement additionally assumes $\chi(g)=1$; the
proof shows that this antecedent is unnecessary for this conclusion.
It remains retained when the statement is returned to its original
consumer. The construction has used the original $AI$ action and
the original embedded $L$, for every supplied $P,AI$.
\end{proof}
% END FET4

\subsection{An integral fixed basis, obtained by finite averages}

For this subsection fix $K,LL,T$. Let $\mathcal N$ be the family
of subgroups $N=H\cap G_L$, where $L/\mathbb Q_p$ is finite
Galois and contains $K$. These are open normal subgroups of $H$,
closed under passing to a common smaller member. The coefficient
action on $W_0$ will be denoted $\sigma_h$, and the semilinear
action on $M_T$ by $\tau_h$.

% BEGIN FET5
\paragraph{FET5: finite quotients act by distinct integral-tilt embeddings.}
\label{claim:FET5}
For $N=H\cap G_L\in\mathcal N$, the maps
\[
 S^N\longrightarrow F,\qquad x\longmapsto h(x),
                   \quad hN\in H/N,                      \tag{FET.4}
\]
are distinct. They have an infinite common set of values obtained
from elements of $S^N$: the powers of $\epsilon-1$, where
$\epsilon$ is the prescribed cyclotomic tilt element.

\begin{proof}
The maps are independent of a chosen representative because $N$
fixes $S^N$. If two maps are equal, $h'^{-1}h$ fixes $S^N$.
Here $N=G_L\cap\ker\chi$, since $L$ contains $K$.
FET4 gives $h'^{-1}h\in G_L$, so it belongs to $N$.
Thus the two cosets are equal.

The prescribed roots defining $\epsilon$ are all fixed by $H$,
so $\epsilon$ and $\epsilon-1$ are fixed. They are the same roots
as in FET1, by the original \texttt{zeta\_eq} field.
The element $\epsilon-1$ is nonzero: a sufficiently high primitive
$p$-power root has $0<v_p(\zeta-1)<1$, and its residue coordinate
is not $1$. It is a nonunit of $S$, because its zeroth residue
coordinate is $0$. In the domain $S$, a nonzero nonunit cannot
have equal distinct nonnegative powers; cancellation would make
one positive power equal to $1$. Their images in $F$ are therefore
an infinite set fixed by every map in (FET.4).
\end{proof}
% END FET5

% BEGIN FET6
\paragraph{FET6: a determinant can be made nonzero by one common average.}
\label{claim:FET6}
Let $A$ be a commutative ring, $C_0$ a field, and
$s_q:A\to C_0$ distinct unital ring maps for a finite nonempty
set $Q$. Suppose an infinite set $B\subset C_0$ consists of
common values: for each $b\in B$ there is $a\in A$ with
$s_q(a)=b$ for every $q$. For square matrices $B_q$ of a fixed
finite size over $C_0$, if $\det B_{q_0}\ne0$ for one $q_0$,
then
\[
              \det\left(\sum_q s_q(y)B_q\right)\ne0
                   \quad\hbox{for some }y\in A.           \tag{FET.5}
\]
The size-zero case is included, with determinant $1$.

\begin{proof}
Distinct ring maps into a field are linearly independent as
functions. For completeness, in a shortest nonzero linear relation
choose two distinct maps, evaluate at $ab$, and subtract one of
their values at $b$ times the relation at $a$. Choosing $b$ where
the two maps differ gives a shorter nonzero relation, a contradiction.
It follows that the vectors $(s_q(a))_q$, as $a$ ranges over $A$,
span $C_0^Q$: a proper span would have a nonzero annihilating
linear functional and hence such a relation.

Choose finitely many $a_1,\ldots,a_m$ whose vectors span, and
consider
\[
 P(X_1,\ldots,X_m)=
 \det\left(\sum_q\sum_{j=1}^m X_j s_q(a_j)B_q\right).
\]
There are coefficients whose vector of inner sums is the unit
vector at $q_0$. At them $P=\det B_{q_0}\ne0$, so $P$ is not
the zero polynomial. A nonzero polynomial over a field cannot
vanish on $B^m$ when $B$ is infinite; prove this by induction
on $m$, using the finite root bound in one variable.
Choose $b_j\in B$ where it does not vanish, and choose common
preimages $c_j\in A$. The element $y=\sum_j c_j a_j$ gives
$s_q(y)=\sum_j b_j s_q(a_j)$ and proves (FET.5).
This proof does not assume that the infinite set is a subfield.
\end{proof}
% END FET6

% BEGIN FET7
\paragraph{FET7: a unit determinant over the original Witt ring.}
\label{claim:FET7}
Let $N\in\mathcal N$, let $Q=H/N$, and choose representatives
$h_q$. For matrices $B_q\in M_d(\mathbb Z_p)$ with one unit
determinant, there exists $y\in W_0^N$ such that
\[
                  \det\left(\sum_q\sigma_{h_q}(y)B_q\right)
                        \in W_0^\times.                  \tag{FET.6}
\]

\begin{proof}
Reduce the matrices modulo $p$ into $F$ and use FET5--FET6
with $A=S^N$. This gives $a\in S^N$ with nonzero determinant
in $F$. Take the Teichmuller lift $y=[a]\in W(F)$.
Functoriality of Witt vectors gives
$\sigma_h([a])=[h(a)]$, so $y$ is $N$-fixed and the determinant
in (FET.6) has the prescribed nonzero zeroth Witt coordinate.
For a perfect field $F$, $W(F)$ is $p$-adically complete and an
element is a unit exactly when its reduction in $F$ is nonzero.
These are the actual Witt-ring properties already established in
WIS2. Thus the determinant is a unit. Neither $|Q|$ nor $p$
has been inverted.
\end{proof}
% END FET7

% BEGIN FET8
\paragraph{FET8: exact trace-one averaging coefficients.}
\label{claim:FET8}
For every $N\in\mathcal N$ there is $b_N\in W_0^N$ with
\[
                      \sum_{q\in H/N}\sigma_{h_q}(b_N)=1.
                                                               \tag{FET.7}
\]

\begin{proof}
Apply FET7 to the one-by-one matrices all equal to $1$.
For the resulting $y$, the sum $u=\sum_q\sigma_{h_q}(y)$ is
a unit. Left multiplication by any $h\in H$ permutes the cosets,
and the right discrepancies of representatives belong to $N$.
Since $y$ is $N$-fixed, this proves $\sigma_h(u)=u$.
The inverse $u^{-1}$ is then fixed as well.
The element $b_N=yu^{-1}$ is $N$-fixed and has trace $1$.
The same coefficient, embedded as $1\otimes b_N$ in
$\mathcal W$, works for $M_T$.
\end{proof}
% END FET8

% BEGIN FET9
\paragraph{FET9: finite averaging improves invariance at a chosen precision.}
\label{claim:FET9}
Let $J_m=p^{m+1}M_T$. Call $x\in M_T$ smooth modulo these
filtrations if for each $m$ some $N\in\mathcal N$ satisfies
$\tau_\nu x-x\in J_m$ for all $\nu\in N$.
If this holds and $N$ has that property at precision $m$, then
for any $c\in\mathcal W^N$ the finite sum
\[
                    A_{N,c}(x)=\sum_{q\in H/N}\tau_{h_q}(cx)
                                                               \tag{FET.8}
\]
is smooth modulo every $J_k$ and is invariant modulo $J_m$
under all of $H$.

\begin{proof}
For $h\in H$, write $hh_q=h_{q'}\nu$ with $\nu\in N$.
Subtract the term indexed by $q'$ after applying $\tau_h$.
Its error is $\tau_{h_{q'}}(c(\tau_\nu x-x))\in J_m$.
Here semilinearity, $\sigma_\nu(c)=c$ and stability of $J_m$
are precisely what is needed. Summing proves invariance modulo
$J_m$.

For smoothness at $J_k$, choose a normal member $N'\subset N$
which fixes $x$ modulo $J_k$. Such a member exists by the
directed property of $\mathcal N$. For $\nu\in N'$, each
conjugate $h_q^{-1}\nu h_q$ lies in $N'$ and in $N$.
It therefore fixes $c$ exactly and $x$ modulo $J_k$.
Subtracting each of the finitely many terms of (FET.8) shows
that $N'$ fixes the average modulo $J_k$. This argument does
not require continuity of arbitrary Witt coefficients.
\end{proof}
% END FET9

% BEGIN FET10
\paragraph{FET10: the averaging limit is exactly invariant.}
\label{claim:FET10}
If $x_0\in M_T$ is smooth modulo the $J_m$, if $N_0$ fixes
it modulo $pM_T$, and if $c_0\in\mathcal W^{N_0}$, then there
exists $x\in M_T^H$ such that
\[
                   x\equiv A_{N_0,c_0}(x_0)\pmod {pM_T}.
                                                               \tag{FET.9}
\]

\begin{proof}
Set $x^{(0)}=A_{N_0,c_0}(x_0)$. FET9 makes it invariant
modulo $J_0$ and smooth. Suppose $x^{(m)}$ is invariant modulo
$J_m$. Choose $N$ fixing it modulo $J_{m+1}$ and average
with the trace-one coefficient in FET8 to define $x^{(m+1)}$.
It is invariant modulo $J_{m+1}$, and
\[
 x^{(m+1)}-x^{(m)}
 =\sum_q\sigma_{h_q}(1\otimes b_N)
                       (\tau_{h_q}x^{(m)}-x^{(m)})\in J_m.
\]
Every iterate is smooth by FET9. The original ring $\mathcal W$
and its finite free module $M_T$ are $p$-adically complete and
separated (WIS2--WIS4). Hence $x^{(m)}$ tends to an $x$ with
$x-x^{(m)}\in J_m$. For each fixed $h$, its semilinear action
preserves every $J_m$. Thus $\tau_hx-x\in J_m$ for every $m$,
and separation makes this difference zero. The congruence at
$m=0$ gives (FET.9).
\end{proof}
% END FET10

% BEGIN FET11
\paragraph{FET11: one initial average is an automorphism.}
\label{claim:FET11}
Choose an $O$-basis $b_1,\ldots,b_r$ of $T$ and a
$\mathbb Z_p$-basis $c_1,\ldots,c_d$ of the same $T$.
There is $N_0\in\mathcal N$ fixing every $b_i$ modulo $pT$.
There is also $y\in W_0^{N_0}$ for which the $\mathcal W$-linear
operator
\[
 U=\sum_{q\in H/N_0}(1\otimes\sigma_{h_q}(y))\rho(h_q)
                         \quad\hbox{on }M_T               \tag{FET.10}
\]
is an automorphism.

\begin{proof}
The lattice action is continuous in its finite
$\mathbb Z_p$-coordinates. Its reduction on $T/pT$ therefore
has an open kernel. An open normal subgroup $G_L$ of
$G_{\mathbb Q_p}$ can be chosen inside the corresponding
neighborhood and inside $G_K$; intersecting with $H$ gives
$N_0$. The same argument works at every $p$-power precision.

In the basis $c$, every $\rho(h_q)$ has an invertible integral
matrix $B_q\in\mathrm{GL}_d(\mathbb Z_p)$. FET7 makes
$\sum_q\sigma_{h_q}(y)B_q$ invertible over $W_0$.
Under the natural identification
$M_T\simeq W_0\otimes_{\mathbb Z_p}T$, this is the matrix of
the underlying $W_0$-linear map $U$. Thus $U$ is bijective.
It commutes with $O$, since all $\rho(h_q)$ do, and with
$W_0$ by its definition. Its inverse therefore commutes with
$\mathcal W$, so it is a $\mathcal W$-automorphism.
This use of a larger $\mathbb Z_p$-basis is only a way to test
invertibility; the eventual basis has the original $O$-rank $r$.
\end{proof}
% END FET11

% BEGIN FET12
\paragraph{FET12: the fixed vectors form a basis on the unchanged lattice.}
\label{claim:FET12}
There is a basis $e_1,\ldots,e_r$ of $M_T$ over $\mathcal W$
fixed by every $h\in H$, where $r=\operatorname{rank}_O T$.

\begin{proof}
For each $1\otimes b_i$, apply FET10 with the common $N_0,y$
of FET11. Smoothness at every precision follows from the
continuity argument there. This gives fixed $e_i$ with
$e_i-U(1\otimes b_i)\in pM_T$.
The columns $U(1\otimes b_i)$ form a $\mathcal W$-basis.
In that basis the matrix of the $e_i$ is $1+pZ$ for an
integral matrix $Z$. The convergent geometric matrix series
$\sum_{n\geq0}(-pZ)^n$ is its inverse.
Thus the $e_i$ are a basis on $M_T$ itself. This argument
uses no division by $p$, no change of $O$ and no enlargement
of $T$. If $r=0$, the empty basis proves the assertion directly.
\end{proof}
% END FET12

% BEGIN FET13
\paragraph{FET13: the full original cyclotomic-kernel fixed-basis theorem.}
\label{claim:FET13}
For every prime $p$, supplied $P,AI$, finite extensions $K,LL$ of $\mathbb Q_p$ and
original lattice $T$, there is a basis indexed by
$\operatorname{Fin}(\operatorname{rank}_{\mathcal O_{LL}}T)$
of the original $\mathrm{TW}(\mathrm{RW}(LL),T)$ satisfying
the original \texttt{FixBasis AI T} predicate.

\begin{proof}
The identifications in (FET.1) are the literal definitions of
$\mathrm{OCb}$, $\mathrm{Cb}$ and $\mathrm{WCb}$. The original
$\mathrm{RW}(LL)$ is $\mathcal W$, and its lattice tensor and
action are exactly $M_T,\tau$. FET1 identifies the given actions;
FET4 supplies faithful invariants for every finite extension.
The family $\mathcal N$ used above is cofinal in the required
identity neighborhoods by the Krull topology and finite Galois
closures. All matrices use the original coefficient embedding.
FET12 therefore gives every quantified field of the exported
\texttt{exists\_cyclotomicKernelFixedBasis\_wittTensorLattice}.
The original implication with antecedent \texttt{FaithH} is
also proved, since FET5 onward only needs that antecedent.
No good-layer, crystalline or Wach-lattice assumption was used
in this fixed-basis theorem.
\end{proof}
% END FET13

\subsection{Constructing the Frobenius conjugacy into the coefficient ring}

Now fix one original good layer $K,s,LL,\Gamma,R,I,\mathrm{IoW}$.
Write $E=R.\mathrm{EE}$, $j=I.\mathrm{ioE}$,
$\mathcal S=\mathcal W^H$, and let $\phi$ be the prescribed
Frobenius. Keep the integral subring
$\mathcal B=O\otimes_{\mathbb Z_p}W(S)$ with its original
inclusion $c:\mathcal B\to\mathcal W$. Put
\[
 t=c(1\otimes([\epsilon_s]-1)),\qquad
 \mathcal N_m=t^m c(\mathcal B)+p\mathcal W\quad(m\geq0).
                                                               \tag{FET.11}
\]
As in IFD1, $t$ is a unit of $\mathcal W$, but (FET.11)
uses $c(\mathcal B)$ and is not the ideal
$t^m\mathcal W+p\mathcal W$. For matrices write
$X\in\mathcal N_m$ when all its entries lie there.

% BEGIN FET14
\paragraph{FET14: the exact two filtrations used in conjugacy.}
\label{claim:FET14}
The actual rings just fixed have the following properties.
The subgroups $\mathcal N_m$ decrease,
$\mathcal N_a\mathcal N_b\subset\mathcal N_{a+b}$,
$\phi(\mathcal N_m)\subset\mathcal N_{pm}$, and
$\bigcap_m\mathcal N_m=p\mathcal W$.
Every finite matrix $X$ over $\mathcal W$ has a common pole bound
$k$ with $t^kX\in\mathcal N_0$.
If $X\in M_r(\mathcal S)$, then for each $m$ there are $N,u,A$
such that
\[
 X-u\in\mathcal N_m,\qquad
                \phi^N(u)-j(A)\in pM_r(\mathcal W),
                         \quad A\in M_r(E).               \tag{FET.12}
\]
The ring $\mathcal S$ is $p$-saturated in $\mathcal W$, closed
under inverses of invertible matrices, and admits sums of sequences
with $i$-th term in $\mathcal S\cap\mathcal N_{i+m}$.
Both $\mathcal S$ and $E$ admit $p$-weighted sums with their
usual congruences, and $\mathcal W$ is $p$-adically separated.

\begin{proof}
IFD1--IFD2 give the multiplication, cancellation and Frobenius
properties of the precise subgroups (FET.11).
The intersection statement follows from WIS23: each zeroth Witt
coordinate belongs to every $\tau_s^m S$, hence is zero by
valuation separation. Thus the element lies in $p\mathcal W$.
Every element of $F$ acquires integral zeroth Witt coordinate
after multiplication by a power of the nonzero $\tau_s$.
Using the fixed finite coefficient basis gives the pole bound for
each entry, and the maximum gives one bound for a finite matrix.

Apply IFD16 to the finitely many entries of $X$, at the same
precision $m$. Choose a common maximum of their Frobenius
iteration exponents. Persistence under Frobenius of $j(E)+p\mathcal W$
(IFD14) gives (FET.12). It is not necessary that $u$ be fixed.
WIS20 gives saturation and $p$-weighted fixed-ring completeness,
and WIS25 gives the stated near-parameter sums at every shift.
If a fixed matrix has an inverse over $\mathcal W$, applying any
$h\in H$ to the inverse equation and using uniqueness shows that
the inverse is fixed. The original $p$-adic completeness of $E$
gives its $p$-weighted sums. WIS2--WIS4 give separation of
$\mathcal W$. These are previously proved properties of the
same $j,t,\phi,H$, not additional analytic assumptions.
\end{proof}
% END FET14

% BEGIN FET15
\paragraph{FET15: the near-parameter contraction solves a matrix equation.}
\label{claim:FET15}
Let $A,B\in M_r(\mathcal S)$ have common pole bound $k$,
and put $m=2k+1$. The additive operator
$\mathcal T(Z)=A\phi(Z)B$ sends
$\mathcal N_{i+m}$ into $\mathcal N_{i+m+1}$.
For $y\in M_r(\mathcal S)\cap\mathcal N_m$ there is
$z\in M_r(\mathcal S)\cap\mathcal N_m$ with
\[
                         z=y+\mathcal T(z)+pR
                      \quad\hbox{for some }R\in M_r(\mathcal W).
                                                               \tag{FET.13}
\]

\begin{proof}
Frobenius raises the filtration index from $i+m$ to $p(i+m)$.
Multiplication by the two matrices loses at most $2k$ by the
cancellation rule in IFD1. Since
$p(i+2k+1)-2k\geq i+2k+2$, the asserted contraction follows.
The sums of $\mathcal T^i(y)$ exist in the fixed ring by FET14.
Let $z$ be such a sum. Its error after $a$ terms lies in
$\mathcal N_{a+m}$, and $z\in\mathcal N_m$ by taking the
empty partial sum. Subtract the finite telescoping identity
from $z-y-\mathcal T(z)$. Contraction shows that this difference
belongs to every $\mathcal N_a$. FET14 puts it in
$pM_r(\mathcal W)$, proving (FET.13).
This is convergence modulo the zeroth Witt coordinate; it is
deliberately separate from the later $p$-adic iteration.
\end{proof}
% END FET15

% BEGIN FET16
\paragraph{FET16: Frobenius iteration makes the approximation sufficiently accurate.}
\label{claim:FET16}
For $F_0\in\mathrm{GL}_r(\mathcal W)$ with entries fixed by
$H$, there are $N\geq0$, $A\in M_r(E)$ and $k\geq0$ such that,
putting $F_1=\phi^N(F_0)$ and $G_1=F_1^{-1}$,
\[
 t^k F_1, t^kG_1, t^k j(A)\in\mathcal N_0,\qquad
                       F_1-j(A)\in\mathcal N_{3k+1}.
                                                               \tag{FET.14}
\]
Also $F_0$ and $F_1$ are Frobenius-conjugate by an invertible
matrix with entries in $\mathcal S$.

\begin{proof}
Choose a common pole bound $k_0$ for $F_0,F_0^{-1}$.
Apply (FET.12) at precision $3k_0+1$, obtaining $N,u,A$.
Frobenius multiplies the error precision by $p^N$.
Since $\phi(t)\equiv t^p\pmod {p\mathcal W}$, applying it
$N$ times to a pole bound multiplies that bound by $p^N$.
Set $k=p^N k_0$. Then
\[
 F_1-j(A)\in\mathcal N_{p^N(3k_0+1)}
                        \subseteq\mathcal N_{3k+1}.
\]
Subtracting this error from $F_1$ gives the same pole bound
for $j(A)$. Invariance and invertibility of $F_1,G_1$ follow
from FET14 and functoriality of Frobenius.

The product
$V_0=F_0\phi(F_0)\cdots\phi^{N-1}(F_0)$, with the empty
product $1$, satisfies
$F_0\phi(V_0)=V_0\phi^N(F_0)$.
It is invertible and fixed, proving the last assertion.
The choice of approximation precision preceded the unknown
iteration exponent $N$; the displayed inequality explains why
this order of choices is valid.
\end{proof}
% END FET16

% BEGIN FET17
\paragraph{FET17: the first conjugacy has coefficient values modulo $p$.}
\label{claim:FET17}
Every invertible fixed matrix $F_0$ is Frobenius-conjugate,
over $\mathcal S$, to an invertible fixed matrix $F_2$ with
\[
                              F_2=j(A)+pR_0
                             \quad(A\in M_r(E)).           \tag{FET.15}
\]
In fact $A$ is invertible over $E$.

\begin{proof}
Use FET16 and put $y=-(F_1-j(A))G_1$.
This is fixed and lies in $\mathcal N_{2k+1}$.
FET15 for the operator $Z\mapsto j(A)\phi(Z)G_1$ gives a
fixed $z\in\mathcal N_{2k+1}$ with
\[
 z=-(F_1-j(A))G_1+j(A)\phi(z)G_1+pR_1.
\]
Multiplying on the right by $F_1$ yields the exact identity
\[
 (1+z)F_1=j(A)\phi(1+z)+pR_1F_1.                           \tag{FET.16}
\]
The matrix $1+z$ is invertible. Indeed the geometric series
in $-z$, interpreted using FET14, gives $W$ with
$(1+z)W=1+pR$. The right side is invertible by the $p$-adic
geometric series in $-pR$. Hence $1+z$ has a right inverse,
and a square matrix over a commutative ring with a right inverse
is invertible (take determinants and the adjugate). Its inverse
$V_1$ is fixed.

Set $F_2=(1+z)F_1\phi(V_1)$. Equation (FET.16) proves
(FET.15), while
$F_1\phi(V_1)=V_1F_2$ gives the required conjugacy.
The matrix $F_2$ is invertible, as a product of invertible
matrices. Then
$j(A)F_2^{-1}=1-pR_0F_2^{-1}$ is invertible, so $j(A)$ is
invertible. WIS18 reflects units through $j$; applied to
$j(\det A)=\det j(A)$, it shows that $A$ is invertible over
the original $E$.
\end{proof}
% END FET17

% BEGIN FET18
\paragraph{FET18: the previously proved quotient theorem supplies every correction.}
\label{claim:FET18}
For $A\in\mathrm{GL}_r(E)$ and every $y\in M_r(\mathcal S)$
there are $z\in M_r(\mathcal S)$ and $d\in M_r(E)$ with
\[
                        y+j(A)\phi(z)-zj(A)=j(d).          \tag{FET.17}
\]

\begin{proof}
Apply the surjectivity in IFD24 to the additive operator
$\Phi(z)=j(A)\phi(z)j(A^{-1})$ on matrices modulo $j(M_r(E))$.
In entry coordinates this is exactly the coefficient matrix
$A\otimes(A^{-1})^{\mathsf t}$ used in IFD25, so IFD24
applies without a new hypothesis. Choose $z$ such that
$\Phi(z)-z$ represents the negative of the class of
$yj(A^{-1})$. Then
$yj(A^{-1})+\Phi(z)-z=j(d_0)$ for some $d_0$ over $E$.
Multiply on the right by $j(A)$ and put $d=d_0A$.
This proves (FET.17) exactly; an error modulo $p$ would already
suffice for the next step.
\end{proof}
% END FET18

% BEGIN FET19
\paragraph{FET19: the integral nonlinear change of basis converges.}
\label{claim:FET19}
For the fixed matrix $F_2=j(A)+pR_0$ in FET17 there are
$V\in\mathrm{GL}_r(\mathcal W)$ with entries in $\mathcal S$
and $A_\infty\in M_r(E)$ such that
\[
 V\equiv1\pmod {p\mathcal W},\qquad
                   F_2\phi(V)=Vj(A_\infty).               \tag{FET.18}
\]

\begin{proof}
Start with $V_0=1$, $D_0=A$. Inductively maintain
$V_n\equiv1\pmod p$, $D_n\equiv A\pmod p$, fixed entries
of $V_n$, and
\[
 F_2\phi(V_n)-V_nj(D_n)=p^{n+1}Y_n.
\]
The left side is fixed. Saturation in FET14, applied repeatedly,
puts every entry of $Y_n$ in $\mathcal S$. Use FET18 to choose
$Z_n,d_n$ with
$Y_n+j(A)\phi(Z_n)-Z_nj(A)=j(d_n)$.
Set
\[
 V_{n+1}=V_n+p^{n+1}Z_n,\qquad
 D_{n+1}=D_n+p^{n+1}d_n.
\]
Expanding the new defect and dividing its expression by
$p^{n+1}$ gives
\[
 Y_n+F_2\phi(Z_n)-V_nj(d_n)-Z_nj(D_n)
                                  -p^{n+1}Z_nj(d_n).
\]
Modulo $p$ this is zero by the choice of $Z_n,d_n$ and the
three maintained congruences. Thus the defect is divisible by
$p^{n+2}$, and all the induction assertions continue.

The $p$-weighted sums in FET14 give limits $V,A_\infty$ in
$M_r(\mathcal S),M_r(E)$ with the stated congruences at every
precision. All ring maps preserve powers of $p$; the finite
matrix operations do as well. Comparing the limiting defect
with the $n$-th defect therefore puts it in every
$p^nM_r(\mathcal W)$. Separation gives zero.
Finally $V=1+pZ$ is invertible by its integral geometric
series. This proves (FET.18) with the fixed coefficients and
without a continuity assertion beyond these explicit congruences.
\end{proof}
% END FET19

% BEGIN FET20
\paragraph{FET20: full original Frobenius-matrix decompletion.}
\label{claim:FET20}
On every original good layer, every
$F_0\in\mathrm{GL}_r(\mathcal W)$ whose entries satisfy
\texttt{FixH} admits
\[
 V\in\mathrm{GL}_r(\mathcal W),\quad A'\in\mathrm{GL}_r(E),
 \qquad V_{ij}\in\mathcal S,\qquad
                          F_0\phi(V)=Vj(A').              \tag{FET.19}
\]
This proves the full original
\texttt{etalePhiMatrix\_decompletion\_to\_wach} endpoint.

\begin{proof}
FET16 conjugates $F_0$ to $F_1$, FET17 conjugates $F_1$ to
$F_2$, and FET19 conjugates $F_2$ to $j(A')$.
If $F\phi(V)=VF'$ and $F'\phi(W)=WF''$, then
$F\phi(VW)=VWF''$. Multiply the three changes of basis in
that order. Each is invertible and fixed, so the resulting $V$
has those properties. Also
$j(A')=V^{-1}F_0\phi(V)$ is invertible. Unit reflection WIS18
again gives $A'\in\mathrm{GL}_r(E)$.
All approximation hypotheses were discharged in FET14 from
WIS/IFD on the actual good layer. No new hypothesis on the
invariant matrix remains. The empty matrix case has $V=A'=1$
in the size-zero matrix ring and satisfies the same conclusion.
\end{proof}
% END FET20

% BEGIN FET21
\paragraph{FET21: full original etale basis on the original tensor lattice.}
\label{claim:FET21}
For every original lattice $T$ on a good layer there is a basis
$e$ of $M_T$ fixed by $H$ whose Frobenius matrix has every
entry in $j(E)$. Thus the original predicate \texttt{EtBasis I T e}
holds at its full scope.

\begin{proof}
Take the fixed basis of FET13. Write
$\Phi(e_j)=\sum_i F_{ij}e_i$, where
$\Phi=\phi\otimes\mathrm{id}_T$.
The Frobenius images span $M_T$: they include
$\Phi(1\otimes t)=1\otimes t$ for every $t\in T$.
Thus the square matrix $F$ has a right inverse and is invertible.
Frobenius commutes with $H$, as one checks on a pure tensor
and on Witt coordinates. The vectors $\Phi(e_j)$ are fixed,
so uniqueness of coordinates in the fixed basis gives
$F_{ij}\in\mathcal S$.

Apply FET20 and define $e'_j=\sum_i V_{ij}e_i$.
The $e'_j$ are a basis and remain fixed. Direct expansion gives
\[
 \Phi(e')=eF\phi(V)=eVj(A')=e'j(A').
\]
Thus the coordinates belong to the image of the original
$I.\mathrm{ioE}$. The matrix orientation is specified in this
calculation; no inverse transpose or changed scalar embedding
has been substituted. This proves every field of \texttt{EtBasis}.
\end{proof}
% END FET21

\subsection{One etale construction, including every lattice morphism}

Choose the basis in FET21 for every lattice $T$, once and for all.
For $r_T=\operatorname{rank}_O T$ define
\[
 D_T=E^{r_T},\qquad
 j_T(x)=\sum_i j(x_i)e_{T,i}\in M_T,
 \qquad \Phi(e_T)=e_Tj(A_T),\quad A_T\in\mathrm{GL}_{r_T}(E).
                                                               \tag{FET.20}
\]
Choice of these bases is permitted by the original existence
statement. The maps below are forced by the embeddings $j_T$, so
they remain coherent despite this choice.

% BEGIN FET22
\paragraph{FET22: the coordinate image and its descent test.}
\label{claim:FET22}
The map $j_T$ is injective, is $E$-semilinear through $j$, and
has $\mathcal W$-spanning image. If a finite list
$x_1,\ldots,x_q\in M_T$ is fixed by $H$ and satisfies
\[
       \Phi(x_j)=\sum_k j(B_{kj})x_k,\qquad B\in\mathrm{GL}_q(E),
                                                               \tag{FET.21}
\]
then every $x_j$ lies in the image of $j_T$.

\begin{proof}
The coordinates of $j_T(x)$ are $j(x_i)$, so injectivity follows
from WIS18. Scalar compatibility is the same coordinate identity.
The image contains all basis vectors $e_{T,i}=j_T(\mathbf e_i)$,
so it spans after extending scalars.
For (FET.21), let $U$ be the matrix of the $x_j$ in the fixed
basis. Its entries are fixed. Taking coordinates of (FET.21)
gives
\[
                       j(A_T)\phi(U)=Uj(B).
\]
IFD25, with zero inhomogeneous term, puts every entry of $U$
in $j(E)$. The corresponding coefficient columns give preimages
under $j_T$. This is descent of an already specified matrix,
which is the distinct role of IFD25 promised at the start.
\end{proof}
% END FET22

% BEGIN FET23
\paragraph{FET23: Frobenius and the quotient Galois action descend.}
\label{claim:FET23}
The image $j_T(D_T)$ is preserved by $\Phi$ and by every
$g\in G_K$. These actions descend to $D_T$ and the latter
action factors through the prescribed group $\Gamma$.
They are semilinear for the prescribed $\phi_E$ and
$\gamma_E$, satisfy the group laws, and commute.

\begin{proof}
The formula for Frobenius in (FET.20), together with
$\phi j=j\phi_E$, gives its descended map directly.
For $g\in G_K$, the list $x_j=g(e_{T,j})$ is fixed by $H$:
the group $H$ is normal and
$h g(e_{T,j})=g(g^{-1}hg)(e_{T,j})=g(e_{T,j})$.
Commutation with Frobenius and the original identity
$g(j(a))=j(\gamma_E(\operatorname{res}_\Gamma g)a)$ give
(FET.21) with
$B=\gamma_E(\operatorname{res}_\Gamma g)(A_T)$, which is
invertible. FET22 descends each basis image, and semilinearity
then descends every vector. Applying the same argument to
$g^{-1}$ gives an inverse.

The original injective cyclotomic character on $\Gamma$ and
its compatibility with $P.\chi_Q$ give
\[
             \operatorname{res}_\Gamma(g)=1
                        \quad\Longleftrightarrow\quad\chi(g)=1.
\]
For the reverse implication use injectivity of that character;
for the forward one use the injective coefficient-field map.
Thus its kernel is precisely $H$, which fixes $j_T(D_T)$
pointwise. Surjectivity of the prescribed
$\operatorname{res}_\Gamma$ gives an action of $\Gamma$:
choose any lift, with independence just proved. Applying $j_T$
reduces the identity, product and commutation laws to those on
$M_T$. Its injectivity returns those laws on $D_T$.
\end{proof}
% END FET23

% BEGIN FET24
\paragraph{FET24: every original lattice morphism descends coherently.}
\label{claim:FET24}
For every original $O$-linear equivariant lattice morphism
$f:T\to T'$, the tensor map $1\otimes f:M_T\to M_{T'}$
induces a unique $E$-linear map $D(f):D_T\to D_{T'}$ with
\[
                       j_{T'}D(f)=(1\otimes f)j_T.        \tag{FET.22}
\]
These maps commute with Frobenius and $\Gamma$, preserve identity,
composition and addition, and send scalar lattice maps to their
prescribed scalar maps.

\begin{proof}
Apply FET22 in $M_{T'}$ to the list
$x_i=(1\otimes f)e_{T,i}$. It is fixed by $H$ because $f$
is equivariant, and satisfies (FET.21) with $B=A_T$ because
the tensor map commutes with Frobenius. Hence it descends.
The prescribed coefficient identity makes the resulting map
$E$-linear, and injectivity of $j_{T'}$ gives uniqueness.

Apply $j_{T'}$ or $j_{T''}$ to each claimed compatibility.
For Frobenius and $\Gamma$ the two compositions become the
commuting actions on the tensor modules. For composition they
both become $(1\otimes(g\circ f))j_T$; for identity, addition
and scalar maps they become the respective elementary tensor
identities. Injectivity proves all equalities. Thus arbitrary
choices of bases have not caused incompatible choices of maps.
\end{proof}
% END FET24

% BEGIN FET25
\paragraph{FET25: every field of the original etale functor.}
\label{claim:FET25}
For every prime $p$, original $P,AI$, and every original good
layer $K,s,LL,\Gamma,R,I,\mathrm{IoW}$, the structure
$\mathrm{WEt}(I)$ is nonempty.

\begin{proof}
Set its module at $T$ to $D_T=E^{r_T}$ from (FET.20).
It is free of the original rank. Its additive group and its
$E$-module structure are the coordinatewise ones; restriction
along the original maps $O\to\mathrm{LetA}\to E$ gives the
two other module structures and both required scalar towers.
FET23 supplies its Frobenius,
quotient Galois action, their semilinearity and all group and
commutation laws. FET24 supplies its map on every original
lattice morphism, its $E$-linearity, the Frobenius and Galois
compatibilities, identity, composition, addition and scalar laws.
The maps $j_T$ in FET22 supply the original injections, scalar
identities and spanning property. Their Frobenius and Galois
intertwining follow from FET23, and their naturality from
(FET.22). These are exactly the fields of the original
\texttt{WEt} record and the full conclusion of
\texttt{EtMod.nonempty\_wEt}.

The only inputs at this endpoint are the original good-layer
records. Faithful invariants, the fixed basis, nonlinear
decompletion and morphism descent have all been proved above or
by the named earlier human proofs. No crystalline representation
is needed to define this etale functor on general lattices.
\end{proof}
% END FET25

% BEGIN FET26
\paragraph{FET26: the actual etale input is now available for the Wach return.}
\label{claim:FET26}
Retain the original crystalline-realization and local-tied-data
premises of WPR30. On every supplied original good layer there
now exist the full original etale functor and its associated
original Wach structure. This does not yet assert the later
Herr or quantitative comparison conclusions.

\begin{proof}
FET25 constructs a $\mathrm{WEt}(I)$ on that layer. WPR30,
which proves the original implication from
\texttt{BhattScholzeCrystallineLattices} and \texttt{LocTied}
to \texttt{WachEx}, applies to this very functor. It gives
the original $\mathrm{NS}$, including its lattice, maps and
comparison fields at the scope already proved in WLD26/WPR30.
The coefficient inclusion, Frobenius, Galois action and lattice
morphisms used by both constructions are the ones just fixed.
The original layerwise quantifier requires this construction for
each allowed layer; it adds no transition-system hypothesis
between distinct outer indices.

The next consumer still invokes the original Herr class maps
and the quantitative comparison. Their existence and bounds do
not follow just from an etale functor or a Wach lattice. Likewise
the common global \texttt{ConclF} family and all of its local
row estimates remain to be supplied before
\texttt{GLOBAL\_DENOMINATOR} can be closed.
\end{proof}
% END FET26

\paragraph{Source dictionary and the next proof.}
The original targets are \path{ConvZero.FaithH},
\path{FixedBasis.main}, \path{MatDecompl.decompletion},
\path{EtaleBasis.exists_etBasis}, and
\path{EtMod.nonempty_wEt}, all on the accepted source pin.
FET14--FET20 prove construction of the nonlinear conjugacy,
whereas IFD25 remains the earlier linear descent theorem used
in FET22. This distinction must remain visible in the roadmap.

For the next step read the actual inputs of
\texttt{ExistsLocalWachFunctor.lean}, then the Herr class maps
and the quantitative local comparison bodies that it calls.
The signed crystalline-lifting estimates must preserve the
full inverse-image finite local condition, including torsion,
the coefficient-normalized length, both cyclic orders of the
signed clusters, and the same original maps. The present result
does not provide those estimates by changing their hypotheses.

\paragraph{Classical foundation references.}
Peter Scholze, \emph{Perfectoid spaces}, Theorem~3.7(i),
\url{https://www.math.uni-bonn.de/people/scholze/PerfectoidSpaces.pdf}.
Bhargav Bhatt, Matthew Morrow and Peter Scholze,
\emph{Integral $p$-adic Hodge theory}, Definition~3.5 and
Example~3.6,
\url{https://www.math.uni-bonn.de/people/scholze/integralpadicHodge.pdf}.
FET3 supplies the elementary sharp-coordinate argument needed
here, so neither a full tilting equivalence nor an unprinted
fixed-field identification is being used as an additional lemma.
